\chapter*{第五部分导读：深度学习、现代 AI 与验证边界}
\addcontentsline{toc}{chapter}{第五部分导读：深度学习、现代 AI 与验证边界}

Part V 进入深度学习和现代 AI。这里的内容看起来会比前面更“前沿”：神经网络、卷积、表示学习、自编码器、Transformer、LLM、agentic workflow、文献阅读辅助和 AI 伦理。但本部分仍然坚持同一个教学原则：先建立可解释 baseline，再引入复杂方法；先说明验证边界，再讨论工具能力。

深度学习在天文和物理中非常有用，因为它能从图像、光谱、序列和复杂特征中学习表示。现代 AI 工具也很有用，因为它们能辅助编程、阅读、整理和检查工作流。但这些能力越强，越需要清楚的验证规则。否则，学生很容易把“模型输出”误当作“科学结论”，把“LLM 草稿”误当作“已经验证的代码或文字”。

\section*{深度学习的概念地图}

本科生第一次接触深度学习时，最容易出现的误解是：把它看成一组模型名字，或者把它看成某个框架里的几行代码。本部分会先把深度学习拆成一组稳定概念，再把这些概念放回天文数据中：

\begin{enumerate}
    \item \textbf{数据表示}：图像、光谱、时间序列和文本都要先变成张量、序列、token 或 patch，模型只能学习它实际看到的表示。
    \item \textbf{网络结构}：全连接层、卷积层、池化层、自编码器瓶颈和注意力层，都是对“哪些信息应该被组合”的结构假设。
    \item \textbf{前向传播}：输入经过一层层变换产生预测、重构或表示。
    \item \textbf{损失函数}：训练并不是让模型“理解科学”，而是让某个可计算目标变小，例如均方误差、交叉熵或重构误差。
    \item \textbf{反向传播与优化}：梯度把误差信号传回参数，优化器根据这些梯度更新权重。
    \item \textbf{泛化与验证}：训练误差下降只说明模型拟合了训练集；验证集、测试集、错误样本和分布外检查才决定结果能否被相信。
    \item \textbf{解释边界}：深度模型学到的是数据中的统计结构和表示，不自动等于物理因果机制。
\end{enumerate}

因此，本部分的深度学习教学目标不是让学生背会所有现代架构，而是让他们能回答五个基本问题：输入是什么形状？模型结构相信了什么？训练目标是什么？模型怎样失败？结果能支持什么科学判断？

\section*{本部分的两条线}

Part V 其实有两条交织的线。

第一条是\textbf{深度表示学习}：从最小神经网络，到图像卷积、一维光谱卷积、自编码器重构、注意力和基础模型概念。它回答的问题是：当手工特征不够时，模型如何学习更有表达力的中间表示？

第二条是\textbf{现代 AI 科研工作流}：LLM 辅助编程、工具调用、文献阅读、报告写作、版权伦理和科研规范。它回答的问题是：当 AI 工具参与科研劳动时，哪些环节可以加速，哪些环节必须由人验证、引用和负责？

\section*{深度模型的学习方式}

本部分不会把深度学习教成一堆框架 API。每一章都应保留四步：

\begin{enumerate}
    \item 先定义一个小型科学任务；
    \item 先做可解释 baseline；
    \item 再引入一个深度学习概念；
    \item 最后回到错误样本、验证集和科学解释。
\end{enumerate}

这能帮助学生理解：神经网络不是魔法，而是一类可训练函数；卷积不是“图像专用咒语”，而是共享局部模式检测；注意力不是万能智能，而是根据上下文动态分配信息权重。

\begin{warnbox}[现代 AI 不能绕过验证]
无论是深度模型还是 LLM 助手，只要它们参与了科研流程，输出就必须经过数据检查、代码运行、引用核验、结果复现和人工解释。速度不是可靠性的替代品。
\end{warnbox}

\section*{本部分的共同产出}

Part V 的每一章都应该留下可复查的记录，而不是只留下“我试过了某个模型”这类模糊印象。本部分不再发明一组新的 chapter record。深度模型继续使用 Part III/IV 已经固定的接口：Dataset Contract、Model Experiment Record、evaluation artifact、error / review artifact 和 Trust Statement。LLM、agent 与 AI 辅助写作则在 Evidence Record 之上增加一份轻量 AI Usage Log，用来说明用了什么材料、由谁核查、哪些输出被接管、哪些内容需要披露。

换句话说，Part V 的目标不是让学生成为模型仓库的使用者，而是成为\textbf{会验证、会披露、会交付}的研究者。
